\documentclass[sigconf]{acmart}
% Cleaning
\pagestyle{fancy} % removes running headers
\settopmatter{printacmref=false} % Removes citation information below abstract
\renewcommand\footnotetextcopyrightpermission[1]{} % removes footnote with conference information in first column
\pagestyle{plain} % removes running headers
\fancyhead{}
\settopmatter{printacmref=false, printfolios=false}
% Copyright
\setcopyright{none}
% \setcopyright{acmcopyright}
% \setcopyright{acmlicensed}
% \setcopyright{rightsretained}
% \setcopyright{usgov}
% \setcopyright{usgovmixed}
% \setcopyright{cagov}
% \setcopyright{cagovmixed}
\settopmatter{printacmref=false} % Removes citation information below abstract
\acmDOI{}


\setlength{\parskip}{0pt}
\setlength{\parsep}{0pt}
\setlength{\headsep}{0pt}
\setlength{\topskip}{0pt}
\setlength{\topmargin}{0pt}
\setlength{\topsep}{0pt}
\setlength{\partopsep}{0pt}
%\linespread{0.95}
\usepackage{mdwlist}
\usepackage{amsmath}
\usepackage{graphicx}
\usepackage{booktabs} % \toprule etc.; harmless if acmart already loads it
\usepackage{array}    %  needed for >{\raggedright\arraybackslash} in the tables
\usepackage{pgfplots}
\pgfplotsset{compat=1.17}


% ---------------------------------------------------------------------------
% [Claude] two helper macros for the draft stage. Both must be gone from the submitted PDF.
%
% \prov{...}  a PROVISIONAL number, printed in red. Every \prov value comes from the current run
%             (evaluation/results/results.json: Oct 16-23 deployment, default settings, seed 0),
%             not from the final test set. Replace each one with the final number, then delete the \prov{ }.
% \note{...}  a note to ourselves, printed in blue. To hide all notes at once, swap the two lines below.
\newcommand{\prov}[1]{\textcolor{red}{#1}}
\newcommand{\note}[1]{{\footnotesize\color{blue}[#1]}}
% \newcommand{\note}[1]{}
% ---------------------------------------------------------------------------
\begin{document}
\title{Evaluating Ocean Drift Prediction from Wind and Location Data}
\author{Jake Cushway}
\affiliation{
  \institution{University of Victoria}
  \city{Victoria}
  \country{Canada}
}
\email{jakecush1@uvic.ca }

\author{Rowan Hall}
\affiliation{
  \institution{University of Victoria}
  \city{Victoria}
  \country{Canada}
}
\email{rowanhall@uvic.ca }


\begin{abstract}
DataXplorer, an uncrewed surface vehicle (USV) serving as a weather station off Half Moon Bay, California, must stay within 10~nautical miles (NM) of the retired NOAA buoy 46012. We forecast its position one hour ahead from onboard wind and GPS data alone, comparing an LSTM and a Transformer with persistence and linear regression on 94~days of 1-minute telemetry. Code and data: \url{https://github.com/jakecush1/autonomous-pilot-ML}.
\end{abstract}

\keywords{Uncrewed surface vehicle, drift prediction, time series regression, LSTM, transformer}
\maketitle
\pagestyle{empty}

\section{Introduction}
Pilots keep DataXplorer on station by watching its position and repositioning it; a reliable drift forecast would let them intervene only when needed. Operational drift models such as GNOME~\cite{zelenke2012gnome} and leeway models~\cite{breivik2011leeway} move an object with a fitted fraction of the wind on top of forecast currents. We instead learn the drift from the USV's own record, and test whether an LSTM~\cite{hochreiter1997lstm} and a Transformer~\cite{vaswani2017attention} improve on persistence and linear regression.


\section{Research Questions}
\noindent\textbf{RQ1.} Do the sequence models predict the USV's position one hour ahead more accurately than the baselines, and which model is best? Answering this shows whether learning patterns within the past hour, such as the drift turning or slowing, is worth the extra complexity. In response, we benchmark the LSTM~\cite{hochreiter1997lstm} and the Transformer~\cite{vaswani2017attention} against persistence and linear regression by their endpoint error on the final test set.


\noindent\textbf{RQ2.} How much does wind contribute to predicting drift? We expected wind to be the main driver of drift, since it both pushes on the USV and drives the surface current.
Leeway models treat it as such~\cite{breivik2011leeway}, and it is one of only two inputs our models use. In response, we measure how often the USV drifts in the same direction as the wind at different wind speeds, how much of the drift a linear fit on the wind can explain, and how the linear regression's error changes without its wind features.




% [Claude] RQ3: "using whole deployments as folds" added (random folds would leak, see the Dataset section).
% Previous sentence: In response, we vary the input window length and learning rate with cross-validation on
% the development set and plot the mean validation error for each setting.
\noindent\textbf{RQ3.} How sensitive are the LSTM and Transformer to their hyperparameters? 

{\color{red}
------------------}

{\color{red}
------------------}

{\color{red}
------------------}

{\color{red}
------------------}

{\color{red}
------------------}

{\color{red}
------------------}



\begin{table}[t]
\centering\footnotesize
\caption{Dataset statistics over all nine deployments.}
\label{tab:dataset}
\begin{tabular}{@{}ll@{}}
\toprule
Deployments            & 9 (Apr 17 -- Nov 16, 2026) \\
Records (1 per minute) & 135{,}234 (93.9 days) \\
% With GPS fix           & 121{,}406 (84.3 days) \\
Missing values         & 11.0\% of records \\
Motor off              & 50.3\% of minutes \\
% Motor-on episodes      & 592 \\
Valid drift minutes    & 66{,}804 (46.4 days) \\
Samples (dev.\ / final test) & 30{,}303 / 3{,}920 \\
% Input features         & 4 per minute $\times$ 60 minutes \\
% Target (next 60 min)   & movement east, north (m) \\
% \quad median / 90th pct / max & 763 / 1{,}469 / 2{,}575 \\
Drift speed (median)   & 0.44\,kn \\
Wind speed (mean / max)& 10.0 / 28.2\,kn \\
Wind direction         & 54\% from the NW \\
\bottomrule
\end{tabular}
\end{table}


\section{Dataset}
The data was collected directly from DataXplorer by Open Ocean Robotics, which operates the vessel at the site of NOAA buoy 46012 off Half Moon Bay, California. We exported only the channels this experiment needs: time, GPS position, wind speed and direction, and motor power. The dataset covers nine deployments between April and November 2026, 135{,}234 one-minute records or 93.9~days on the water over seven months. Since only drift is of interest, a minute is kept only when the motor is off, no value is missing and the vessel moves slower than 3~kn (faster means a tow). A sample needs 120 consecutive such minutes: the past hour is the input, with four features per minute (movement and wind push, each east and north), and the movement over the next hour is the target, read from later GPS positions without labelling. This yields 34{,}223 samples. From Table~\ref{tab:dataset}, the vessel drifts half of the time, and the hourly drift has a long tail, with a median of 763~m and a maximum of 2{,}575~m. No public benchmark covers this vessel, but its own record at the site where the forecast will be used is the most relevant data; it spans spring to autumn and winds from calm to 28~kn, though it covers one vessel, no winter, and no measurement of current.

% ===========================================================================
% The headings follow the manual (4.4 Baselines, 4.5 Evaluation Strategy and Metrics, 4.6 Results).
% Marks: baselines 5%, strategy and metrics 5%, results with attribution and discussion 50%, reproducibility 20%.
% ===========================================================================


\section{Baselines}

% NOTE (manual 4.4): for each baseline say how it works, why it is a fair competitor, its hyperparameters,
% the library, and a reference. Milestone 3 has the full descriptions, so a sentence or two each is enough.
We compare the LSTM and Transformer with two baselines. All four models use the same samples, cleaning and metric.

\textbf{Persistence} predicts that the next hour's movement repeats the last hour's. It needs no training and has no hyperparameters. It is the bar every model must beat.

\textbf{Linear regression} predicts the next hour's movement from six numbers summarising the past hour: the movement, the current wind and the mean wind. Like a leeway model~\cite{breivik2011leeway}, it moves the USV linearly with the wind, on top of a term for its past drift, so it tests whether a linear model suffices before turning to nonlinear ones. We use scikit-learn's \texttt{LinearRegression}, which has no hyperparameters to tune.


\textbf{LSTM and Transformer.}
% [Claude] NOTE for Rowan (this part is yours to write). Facts from src/models/nn_models.py:
%   both read the past hour minute by minute ([dx, dy, wind_push_e, wind_push_n] per minute);
%   LSTM: 2 layers, 64 hidden, dropout 0.2, last hidden state -> (east, north);
%   Transformer: 2 encoder layers, d=64, 4 heads, sin/cos positions, mean-pool -> (east, north);
%   PyTorch, Adam, MSE loss, inputs/targets standardised with train stats, seed 0.
%   Possible draft, if useful: "Both read the past hour minute by minute, so they can see the drift turn or slow
%   within the hour. The LSTM~\cite{hochreiter1997lstm} carries a memory through the sequence, while the
%   Transformer~\cite{vaswani2017attention} uses attention to weigh every minute against every other. Each has
%   two layers of width 64, is built in PyTorch, and is trained with Adam~\cite{kingma2015adam} on the mean squared error."
% Both read the past hour minute by minute. Each has 2 layers of width 64 in PyTorch and is trained with Adam~\cite{kingma2015adam} and early stopping. Their input window and learning rate are studied in RQ3.

{\color{red}
------------------}


{\color{red}
------------------}


{\color{red}
------------------}


% [Claude] Table 2 replaced by a bar chart (final test set, evaluation/results/results.json "test").
% The old table (Oct 16-23 numbers) is in git history.
\begin{figure}[t]
\centering
\begin{tikzpicture}
\begin{axis}[
    xbar, bar width=6pt,
    scale only axis, width=0.48\columnwidth, height=4.2cm,  % axis box only; the model names add ~3.3 cm
    xmin=0, xmax=700, xtick={0,200,400,600},
    xlabel={Endpoint error after 60 minutes (m)},
    ytick={1,2,3,4,5},
    yticklabels={{Persistence (0.87)}, {Linear reg.\ no wind (0.86)}, {Linear regression (0.90)}, {LSTM (0.92)}, {Transformer (0.93)}},
    yticklabel style={font=\footnotesize},
    y dir=reverse,
    enlarge y limits=0.15,
    axis lines*=left,
    nodes near coords,
    nodes near coords style={font=\scriptsize},
    legend style={at={(0.5,1.02)}, anchor=south, legend columns=2, font=\scriptsize, draw=none},
]
\addplot[fill=gray!30, draw=black] coordinates {(305,1) (301,2) (269,3) (213,4) (222,5)};
\addplot[fill=gray!70, draw=black] coordinates {(572,1) (587,2) (476,3) (433,4) (387,5)};
\legend{median, 90th percentile}
\end{axis}
\end{tikzpicture}
\Description{Horizontal bar chart of the median and 90th-percentile endpoint error on the final test set:
persistence 305 and 572 m (R-squared 0.87), linear regression without wind 301 and 587 m (0.86), linear regression
269 and 476 m (0.90), LSTM 213 and 433 m (0.92), Transformer 222 and 387 m (0.93).}
\caption{Endpoint error after 60 minutes on the final test set (November), with $R^2$ in brackets. Assuming no drift gives 747~m (median) and 1{,}663~m (90th percentile).}
\label{fig:results}
\end{figure}


\section{Evaluation Strategy and Metrics}

\textbf{Splits.} Samples one minute apart share 59 of their 60 input minutes, so a random split would place near-copies of test samples in training. We therefore split by whole deployment in time order: the seven deployments from April to October form the development set (30{,}303 samples, 89\%) and the two from November the final test set (3{,}920 samples, 11\%). Every choice is made on the development set, and the final test set is scored once, after the settings are fixed.

\textbf{Metric.} We score a prediction by its endpoint error, the distance in metres between the predicted and the true position after 60 minutes. We report the median as the typical miss and the 90th percentile as a bad case that a pilot could use as a safety radius. We also report $R^2$, the share of the variation in the movement east and north that a model explains (1 is perfect, 0 is no better than always predicting the average movement), as the standard measure for regression. The output is a position, so this is regression, and classification metrics such as accuracy or F1 do not apply.

\textbf{Hyperparameters.}
% [Claude] NOTE for Rowan (this part is yours to write). Facts for the final run (Nov 18):
%   learning rate 1e-3, batch 256, up to 80 epochs, early stopping with patience 10 on the validation
%   deployment Oct 16-23 (changed from Oct 3-13: train is now Apr-Oct 13, test = November holdout);
%   no hyperparameter search was run, so these are untuned defaults; the tab:hyper table is commented out below.
%   Persistence and linear regression have no hyperparameters.

{\color{red}
------------------}

{\color{red}
------------------}

{\color{red}
------------------}

% \begin{table}[t]
% \centering\footnotesize
% % NOTE: mark the best value of each row in bold once RQ3 is run.
% \caption{Hyperparameters of the LSTM and Transformer (RQ3).}
% \label{tab:hyper}
% \begin{tabular}{@{}ll@{}}
% \toprule
% Hyperparameter & Values \\
% \midrule
% Input window (min) & \prov{15, 30, 60, 120} \\
% Learning rate & \prov{$10^{-4}$, $3{\times}10^{-4}$, $10^{-3}$, $3{\times}10^{-3}$} \\
% \bottomrule
% \end{tabular}
% \end{table}

\section{Results and Discussion}



% NOTE (manual 4.6): one part for each RQ: state the answer, give the numbers, then say why.
% No result that does not answer an RQ. This section is 50% of the mark, so it gets the most space.

% [Claude] Nov 18 2026: RQ1 now on the FINAL TEST SET (data/holdout, Nov 2-16), scored once; train Apr-Oct 13,
% early stopping on Oct 16-23, seed 0. Numbers: evaluation/results/results.json ("test"). The old red Oct 16-23
% numbers are in git history. Still one seed per model, so the 9 m LSTM-Transformer gap may be noise.

% [Claude] RQ1 answer with the final-test numbers, in the example paper's pattern.
\textbf{RQ1: do the sequence models beat the baselines?} We summarize the results on the final test set in Figure~\ref{fig:results}. As shown, in terms of median error, both sequence models outperform all baselines, and the LSTM is best at 213~m, 30\% below persistence and 21\% below linear regression. In terms of the 90th percentile, however, the Transformer is best at 387~m, 32\% below persistence. In terms of $R^2$, the sequence models explain 92--93\% of the variation in the movement, against 87\% for persistence and 90\% for linear regression. This can be attributed to the sequence models reading the past hour minute by minute: they see the drift turning or slowing, whereas persistence assumes the last hour repeats and linear regression sees only its totals. Persistence alone more than halves the error of assuming no drift (747~m), because the current, which drives most of the drift (RQ2), changes slowly. Linear regression also beats persistence (269~m), but only through its wind features, as RQ2 shows. Overall, the LSTM is the best model by median error and the Transformer by the 90th percentile; with one training seed per model, the 9~m gap between them may be noise.



% [Claude] RQ2 answer, in the manual's 4.6 pattern (answer, numbers, then why). Numbers from
% src/linear_fits/drift_vs_wind.py 60 (development set, 25,640 motor-off minutes, drift over 60 min).
% Earlier placeholder, kept in case there is room for the run without wind:
% NOTE: show the error with and without wind for each learned model, as a small table or a bar chart.
% \textbf{RQ2: how much does wind contribute?} We retrain each learned model with the two wind features removed. \prov{Without wind, the median error of the [model] changes from [A]~m to [B]~m.} \prov{[Why.]}
\noindent\textbf{RQ2. how much does wind contribute?} In response to RQ2, Figure~\ref{fig:withwind} shows that wind steers the drift but does not dominate it. Over the development set, with the drift measured over 60 minutes, the USV drifts with the wind (within $\pm45^\circ$ of downwind) 47\% of the time, against the 25\% expected if wind had no effect on direction. The effect grows with wind speed: below 5~kn the share is 20\%, no better than chance, while above 15~kn it reaches 87\%. This can be attributed to the balance between the wind's push and the current: in light wind the current sets the direction, and only a strong wind overcomes it. For the same reason, wind speed alone explains little of the drift speed ($R^2 = 0.11$), because a current in any direction adds to the speed. The part of the drift along the wind grows by 4.2\% of the wind speed, just above the 1--4\% windage that GNOME assumes~\cite{zelenke2012gnome}, and a fit on the wind vector explains more of the drift ($R^2 = 0.23$), turning it about $20^\circ$ to the right of downwind, as expected from the Coriolis effect in the Northern Hemisphere. Wind thus accounts for about a quarter of the drift and current and tide for the rest, which explains why past drift, which already carries the current, is the strongest input. Its value for prediction therefore depends on the weather. On the final test set, removing the wind features from the linear regression raises its median error from 269 to 301~m and its 90th percentile from 476 to 587~m, and lowers its $R^2$ from 0.90 to 0.86, so that it is no better than persistence, whereas on the calmer validation week they change its median by only 2~m. This supports our explanation: the November deployments were much windier, with an hourly mean wind of at least 10~kn in 73\% of samples against 34\% in the validation week.

\begin{figure}[t]
\centering
\begin{tikzpicture}
\begin{axis}[
    ybar, bar width=14pt,
    scale only axis, width=0.78\columnwidth, height=3.6cm,  % axis box only; the y label and ticks add ~1.3 cm
    ymin=0, ymax=100, ytick={0,25,50,75,100},
    ylabel={Time drifting with the wind (\%)},
    xlabel={Wind speed (kn)},
    xtick={1,2,3,4,5},
    xticklabels={{1--5\\\scriptsize(19\%)}, {5--10\\\scriptsize(55\%)},
                 {10--15\\\scriptsize(23\%)}, {$\geq$15\\\scriptsize(3\%)}, {All\\\scriptsize(100\%)}},
    xticklabel style={align=center, font=\footnotesize},
    enlarge x limits=0.15,
    axis lines*=left,
    nodes near coords={\pgfmathprintnumber{\pgfplotspointmeta}\%},
    nodes near coords style={font=\footnotesize},
    legend style={at={(0.5,1.03)}, anchor=south, legend columns=2, font=\scriptsize, draw=none},
]
\addplot[fill=gray!55, draw=black] coordinates {(1,20) (2,45) (3,64) (4,87) (5,47)};
\addlegendentry{drift within $\pm45^\circ$ of downwind}
\addlegendimage{line legend, dashed, black}
\addlegendentry{no wind effect (25\%)}
\draw[dashed] (rel axis cs:0,0.25) -- (rel axis cs:1,0.25);
\end{axis}
\end{tikzpicture}
\Description{Bar chart of the share of time the USV drifts within 45 degrees of downwind:
20\% at 1--5 kn, 45\% at 5--10 kn, 64\% at 10--15 kn, 87\% above 15 kn and 47\% overall,
against 25\% expected if wind had no effect on direction.}
\caption{Time the USV drifts within $\pm45^\circ$ of downwind, by wind speed
(share of minutes in brackets). Dashed: no wind effect.}
\label{fig:withwind}
\end{figure}

% NOTE: one figure: mean cross-validation error against the input window and against the learning rate,
% one line for each model. Say which setting is best, how flat the curve is, and why.
% [Claude] NOTE for Rowan: fig:hyper no longer exists, so this \ref prints "Figure ??" until the figure is added
% or the sentence changes. No sweep has been run yet (see the Hyperparameters note).
\textbf{RQ3: how sensitive are the models?} Figure~\ref{fig:hyper} shows the mean cross-validation error for each setting. 

{\color{red}
------------------}

{\color{red}
------------------}

{\color{red}
------------------}



% NOTE: 20% of the mark. Give the one command that reproduces every table and figure, the fixed seeds,
% the run time, and where the results are written.
\textbf{Reproducibility.} The command \texttt{bash quickstart.sh} installs the pinned dependencies and reruns every step, from the raw telemetry in \texttt{data/} to the scores in \texttt{evaluation/results/} and the trained models in \texttt{output/}. Seeds are fixed, and the full run takes about 5~minutes on a laptop CPU.

% NOTE: if there is room, close with one sentence: what a pilot can rely on, and what comes next.

\bibliographystyle{ACM-Reference-Format}
\bibliography{bibliography.bib}

\end{document}
\documentclass[sigconf]{acmart}
% Cleaning
\pagestyle{fancy} % removes running headers
\settopmatter{printacmref=false} % Removes citation information below abstract
\renewcommand\footnotetextcopyrightpermission[1]{} % removes footnote with conference information in first column
\pagestyle{plain} % removes running headers
\fancyhead{}
\settopmatter{printacmref=false, printfolios=false}
% Copyright
\setcopyright{none}
% \setcopyright{acmcopyright}
% \setcopyright{acmlicensed}
% \setcopyright{rightsretained}
% \setcopyright{usgov}
% \setcopyright{usgovmixed}
% \setcopyright{cagov}
% \setcopyright{cagovmixed}
\settopmatter{printacmref=false} % Removes citation information below abstract
\acmDOI{}


\setlength{\parskip}{0pt}
\setlength{\parsep}{0pt}
\setlength{\headsep}{0pt}
\setlength{\topskip}{0pt}
\setlength{\topmargin}{0pt}
\setlength{\topsep}{0pt}
\setlength{\partopsep}{0pt}
%\linespread{0.95}
\usepackage{mdwlist}
\usepackage{amsmath}
\usepackage{graphicx}
\usepackage{booktabs} % \toprule etc.; harmless if acmart already loads it
\usepackage{array}    %  needed for >{\raggedright\arraybackslash} in the tables
\usepackage{pgfplots}
\pgfplotsset{compat=1.17}


% ---------------------------------------------------------------------------
% [Claude] two helper macros for the draft stage. Both must be gone from the submitted PDF.
%
% \prov{...}  a PROVISIONAL number, printed in red. Every \prov value comes from the current run
%             (evaluation/results/results.json: Oct 16-23 deployment, default settings, seed 0),
%             not from the final test set. Replace each one with the final number, then delete the \prov{ }.
% \note{...}  a note to ourselves, printed in blue. To hide all notes at once, swap the two lines below.
\newcommand{\prov}[1]{\textcolor{red}{#1}}
\newcommand{\note}[1]{{\footnotesize\color{blue}[#1]}}
% \newcommand{\note}[1]{}
% ---------------------------------------------------------------------------
\begin{document}
\title{Evaluating Ocean Drift Prediction from Wind and Location Data}
\author{Jake Cushway}
\affiliation{
  \institution{University of Victoria}
  \city{Victoria}
  \country{Canada}
}
\email{jakecush1@uvic.ca }

\author{Rowan Hall}
\affiliation{
  \institution{University of Victoria}
  \city{Victoria}
  \country{Canada}
}
\email{rowanhall@uvic.ca }


\begin{abstract}
DataXplorer, an uncrewed surface vehicle (USV) serving as a weather station off Half Moon Bay, California, must stay within 10~nautical miles (NM) of the retired NOAA buoy 46012. We forecast its position one hour ahead from onboard wind and GPS data alone, comparing an LSTM and a Transformer with persistence and linear regression on 94~days of 1-minute telemetry. Code and data: \url{https://github.com/jakecush1/autonomous-pilot-ML}.
\end{abstract}

\keywords{Uncrewed surface vehicle, drift prediction, time series regression, LSTM, transformer}
\maketitle
\pagestyle{empty}

\section{Introduction}
Pilots keep DataXplorer on station by watching its position and repositioning it; a reliable drift forecast would let them intervene only when needed. Operational drift models such as GNOME~\cite{zelenke2012gnome} and leeway models~\cite{breivik2011leeway} move an object with a fitted fraction of the wind on top of forecast currents. We instead learn the drift from the USV's own record, and test whether an LSTM~\cite{hochreiter1997lstm} and a Transformer~\cite{vaswani2017attention} improve on persistence and linear regression.


\section{Research Questions}
\noindent\textbf{RQ1.} Do the sequence models predict the USV's position one hour ahead more accurately than the baselines, and which model is best? Answering this shows whether learning patterns within the past hour, such as the drift turning or slowing, is worth the extra complexity. In response, we benchmark the LSTM~\cite{hochreiter1997lstm} and the Transformer~\cite{vaswani2017attention} against persistence and linear regression by their endpoint error on the final test set.


\noindent\textbf{RQ2.} How much does wind contribute to predicting drift? We expected wind to be the main driver of drift, since it both pushes on the USV and drives the surface current.
Leeway models treat it as such~\cite{breivik2011leeway}, and it is one of only two inputs our models use. In response, we measure how often the USV drifts in the same direction as the wind at different wind speeds, how much of the drift a linear fit on the wind can explain, and how the linear regression's error changes without its wind features.




% [Claude] RQ3: "using whole deployments as folds" added (random folds would leak, see the Dataset section).
% Previous sentence: In response, we vary the input window length and learning rate with cross-validation on
% the development set and plot the mean validation error for each setting.
\noindent\textbf{RQ3.} How sensitive are the LSTM and Transformer to their hyperparameters? 

{\color{red}
------------------}

{\color{red}
------------------}

{\color{red}
------------------}

\noindent\textbf{RQ4.} How is this method different from widely used methods? 

% Other popular methods will use large predictive current models, and current forecasts to predict the future drift.  Our's is a much lighter approach and uses only data which is readily available onboard the USV, namely, wind speed, wind direction, and location.

{\color{red}
------------------}

{\color{red}
------------------}


{\color{red}
------------------}


{\color{red}
------------------}


\begin{table}[t]
\centering\footnotesize
\caption{Dataset statistics over all nine deployments.}
\label{tab:dataset}
\begin{tabular}{@{}ll@{}}
\toprule
Deployments            & 9 (Apr 17 -- Nov 16, 2026) \\
Records (1 per minute) & 135{,}234 (93.9 days) \\
% With GPS fix           & 121{,}406 (84.3 days) \\
Missing values         & 11.0\% of records \\
Motor off              & 50.3\% of minutes \\
% Motor-on episodes      & 592 \\
Valid drift minutes    & 66{,}804 (46.4 days) \\
Samples (dev.\ / final test) & 30{,}303 / 3{,}920 \\
% Input features         & 4 per minute $\times$ 60 minutes \\
% Target (next 60 min)   & movement east, north (m) \\
% \quad median / 90th pct / max & 763 / 1{,}469 / 2{,}575 \\
Drift speed (median)   & 0.44\,kn \\
Wind speed (mean / max)& 10.0 / 28.2\,kn \\
Wind direction         & 54\% from the NW \\
\bottomrule
\end{tabular}
\end{table}


\section{Dataset}
The data was collected directly from DataXplorer by Open Ocean Robotics, which operates the vessel at the site of NOAA buoy 46012 off Half Moon Bay, California. We exported only the channels this experiment needs: time, GPS position, wind speed and direction, and motor power. The dataset covers nine deployments between April and November 2026, 135{,}234 one-minute records or 93.9~days on the water over seven months. Since only drift is of interest, a minute is kept only when the motor is off, no value is missing and the vessel moves slower than 3~kn (faster means a tow). A sample needs 120 consecutive such minutes: the past hour is the input, with four features per minute (movement and wind push, each east and north), and the movement over the next hour is the target, read from later GPS positions without labelling. This yields 34{,}223 samples. From Table~\ref{tab:dataset}, the vessel drifts half of the time, and the hourly drift has a long tail, with a median of 763~m and a maximum of 2{,}575~m. No public benchmark covers this vessel, but its own record at the site where the forecast will be used is the most relevant data; it spans spring to autumn and winds from calm to 28~kn, though it covers one vessel, no winter, and no measurement of current.

% ===========================================================================
% The headings follow the manual (4.4 Baselines, 4.5 Evaluation Strategy and Metrics, 4.6 Results).
% Marks: baselines 5%, strategy and metrics 5%, results with attribution and discussion 50%, reproducibility 20%.
% ===========================================================================


\section{Baselines}

% NOTE (manual 4.4): for each baseline say how it works, why it is a fair competitor, its hyperparameters,
% the library, and a reference. Milestone 3 has the full descriptions, so a sentence or two each is enough.
We compare the LSTM and Transformer with two baselines. All four models use the same samples, cleaning and metric.

\textbf{Persistence} predicts that the next hour's movement repeats the last hour's. It needs no training and has no hyperparameters. It is the bar every model must beat.

\textbf{Linear regression} predicts the next hour's movement from six numbers summarising the past hour: the movement, the current wind and the mean wind. Like a leeway model~\cite{breivik2011leeway}, it moves the USV linearly with the wind, on top of a term for its past drift, so it tests whether a linear model suffices before turning to nonlinear ones. We use scikit-learn's \texttt{LinearRegression}, which has no hyperparameters to tune.


\textbf{LSTM and Transformer.}
% [Claude] NOTE for Rowan (this part is yours to write). Facts from src/models/nn_models.py:
%   both read the past hour minute by minute ([dx, dy, wind_push_e, wind_push_n] per minute);
%   LSTM: 2 layers, 64 hidden, dropout 0.2, last hidden state -> (east, north);
%   Transformer: 2 encoder layers, d=64, 4 heads, sin/cos positions, mean-pool -> (east, north);
%   PyTorch, Adam, MSE loss, inputs/targets standardised with train stats, seed 0.
%   Possible draft, if useful: "Both read the past hour minute by minute, so they can see the drift turn or slow
%   within the hour. The LSTM~\cite{hochreiter1997lstm} carries a memory through the sequence, while the
%   Transformer~\cite{vaswani2017attention} uses attention to weigh every minute against every other. Each has
%   two layers of width 64, is built in PyTorch, and is trained with Adam~\cite{kingma2015adam} on the mean squared error."
% Both read the past hour minute by minute. Each has 2 layers of width 64 in PyTorch and is trained with Adam~\cite{kingma2015adam} and early stopping. Their input window and learning rate are studied in RQ3.

{\color{red}
------------------}


{\color{red}
------------------}


{\color{red}
------------------}

\section{Evaluation Strategy and Metrics}

\textbf{Splits.} Samples one minute apart share 59 of their 60 input minutes, so a random split would place near-copies of test samples in training. We therefore split by whole deployment in time order: the seven deployments from April to October form the development set (30{,}303 samples, 89\%) and the two from November the final test set (3{,}920 samples, 11\%). Every choice is made on the development set, and the final test set is scored once, after the settings are fixed.

\textbf{Metric.} We score a prediction by its endpoint error, the distance in metres between the predicted and the true position after 60 minutes. We report the median as the typical miss and the 90th percentile as a bad case that a pilot could use as a safety radius. We also report $R^2$, the share of the variation in the movement east and north that a model explains (1 is perfect, 0 is no better than always predicting the average movement), as the standard measure for regression. The output is a position, so this is regression, and classification metrics such as accuracy or F1 do not apply.

\textbf{Hyperparameters.}
% [Claude] NOTE for Rowan (this part is yours to write). Facts for the final run (Nov 18):
%   learning rate 1e-3, batch 256, up to 80 epochs, early stopping with patience 10 on the validation
%   deployment Oct 16-23 (changed from Oct 3-13: train is now Apr-Oct 13, test = November holdout);
%   no hyperparameter search was run, so these are untuned defaults; the tab:hyper table is commented out below.
%   Persistence and linear regression have no hyperparameters.

{\color{red}
------------------}

{\color{red}
------------------}

{\color{red}
------------------}

% \begin{table}[t]
% \centering\footnotesize
% % NOTE: mark the best value of each row in bold once RQ3 is run.
% \caption{Hyperparameters of the LSTM and Transformer (RQ3).}
% \label{tab:hyper}
% \begin{tabular}{@{}ll@{}}
% \toprule
% Hyperparameter & Values \\
% \midrule
% Input window (min) & \prov{15, 30, 60, 120} \\
% Learning rate & \prov{$10^{-4}$, $3{\times}10^{-4}$, $10^{-3}$, $3{\times}10^{-3}$} \\
% \bottomrule
% \end{tabular}
% \end{table}

\section{Results and Discussion}



% NOTE (manual 4.6): one part for each RQ: state the answer, give the numbers, then say why.
% No result that does not answer an RQ. This section is 50% of the mark, so it gets the most space.

% [Claude] Nov 18 2026: RQ1 now on the FINAL TEST SET (data/holdout, Nov 2-16), scored once; train Apr-Oct 13,
% early stopping on Oct 16-23, seed 0. Numbers: evaluation/results/results.json ("test"). The old red Oct 16-23
% numbers are in git history. Still one seed per model, so the 9 m LSTM-Transformer gap may be noise.

% [Claude] RQ1 answer with the final-test numbers, in the example paper's pattern.
\textbf{RQ1: do the sequence models beat the baselines?} We summarize the results on the final test set in Figure~\ref{fig:results}. As shown, in terms of median error, both sequence models outperform all baselines, and the LSTM is best at 213~m, 30\% below persistence and 21\% below linear regression. In terms of the 90th percentile, however, the Transformer is best at 387~m, 32\% below persistence. In terms of $R^2$, the sequence models explain 92--93\% of the variation in the movement, against 87\% for persistence and 90\% for linear regression. This can be attributed to the sequence models reading the past hour minute by minute: they see the drift turning or slowing, whereas persistence assumes the last hour repeats and linear regression sees only its totals. Persistence alone more than halves the error of assuming no drift (747~m), because the current, which drives most of the drift (RQ2), changes slowly. Linear regression also beats persistence (269~m), but only through its wind features, as RQ2 shows. Overall, the LSTM is the best model by median error and the Transformer by the 90th percentile; with one training seed per model, the 9~m gap between them may be noise.


% [Claude] Table 2 replaced by a bar chart (final test set, evaluation/results/results.json "test").
% The old table (Oct 16-23 numbers) is in git history.
\begin{figure}[t]
\centering
\begin{tikzpicture}
\begin{axis}[
    xbar, bar width=6pt,
    scale only axis, width=0.48\columnwidth, height=4.2cm,  % axis box only; the model names add ~3.3 cm
    xmin=0, xmax=700, xtick={0,200,400,600},
    xlabel={Endpoint error after 60 minutes (m)},
    ytick={1,2,3,4,5},
    yticklabels={{Persistence (0.87)}, {Linear reg.\ no wind (0.86)}, {Linear regression (0.90)}, {LSTM (0.92)}, {Transformer (0.93)}},
    yticklabel style={font=\footnotesize},
    y dir=reverse,
    enlarge y limits=0.15,
    axis lines*=left,
    nodes near coords,
    nodes near coords style={font=\scriptsize},
    legend style={at={(0.5,1.02)}, anchor=south, legend columns=2, font=\scriptsize, draw=none},
]
\addplot[fill=gray!30, draw=black] coordinates {(305,1) (301,2) (269,3) (213,4) (222,5)};
\addplot[fill=gray!70, draw=black] coordinates {(572,1) (587,2) (476,3) (433,4) (387,5)};
\legend{median, 90th percentile}
\end{axis}
\end{tikzpicture}
\Description{Horizontal bar chart of the median and 90th-percentile endpoint error on the final test set:
persistence 305 and 572 m (R-squared 0.87), linear regression without wind 301 and 587 m (0.86), linear regression
269 and 476 m (0.90), LSTM 213 and 433 m (0.92), Transformer 222 and 387 m (0.93).}
\caption{Endpoint error after 60 minutes on the final test set (November), with $R^2$ in brackets. Assuming no drift gives 747~m (median) and 1{,}663~m (90th percentile).}
\label{fig:results}
\end{figure}

% [Claude] RQ2 answer, in the manual's 4.6 pattern (answer, numbers, then why). Numbers from
% src/linear_fits/drift_vs_wind.py 60 (development set, 25,640 motor-off minutes, drift over 60 min).
% Earlier placeholder, kept in case there is room for the run without wind:
% NOTE: show the error with and without wind for each learned model, as a small table or a bar chart.
% \textbf{RQ2: how much does wind contribute?} We retrain each learned model with the two wind features removed. \prov{Without wind, the median error of the [model] changes from [A]~m to [B]~m.} \prov{[Why.]}
\noindent\textbf{RQ2. how much does wind contribute?} In response to RQ2, Figure~\ref{fig:withwind} shows that wind steers the drift but does not dominate it. Over the development set, with the drift measured over 60 minutes, the USV drifts with the wind (within $\pm45^\circ$ of downwind) 47\% of the time, against the 25\% expected if wind had no effect on direction. The effect grows with wind speed: below 5~kn the share is 20\%, no better than chance, while above 15~kn it reaches 87\%. This can be attributed to the balance between the wind's push and the current: in light wind the current sets the direction, and only a strong wind overcomes it. For the same reason, wind speed alone explains little of the drift speed ($R^2 = 0.11$), because a current in any direction adds to the speed. The part of the drift along the wind grows by 4.2\% of the wind speed, just above the 1--4\% windage that GNOME assumes~\cite{zelenke2012gnome}, and a fit on the wind vector explains more of the drift ($R^2 = 0.23$), turning it about $20^\circ$ to the right of downwind, as expected from the Coriolis effect in the Northern Hemisphere. Wind thus accounts for about a quarter of the drift and current and tide for the rest, which explains why past drift, which already carries the current, is the strongest input. Its value for prediction therefore depends on the weather. On the final test set, removing the wind features from the linear regression raises its median error from 269 to 301~m and its 90th percentile from 476 to 587~m, and lowers its $R^2$ from 0.90 to 0.86, so that it is no better than persistence, whereas on the calmer validation week they change its median by only 2~m. This supports our explanation: the November deployments were much windier, with an hourly mean wind of at least 10~kn in 73\% of samples against 34\% in the validation week.

\begin{figure}[t]
\centering
\begin{tikzpicture}
\begin{axis}[
    ybar, bar width=14pt,
    scale only axis, width=0.78\columnwidth, height=3.6cm,  % axis box only; the y label and ticks add ~1.3 cm
    ymin=0, ymax=100, ytick={0,25,50,75,100},
    ylabel={Time drifting with the wind (\%)},
    xlabel={Wind speed (kn)},
    xtick={1,2,3,4,5},
    xticklabels={{1--5\\\scriptsize(19\%)}, {5--10\\\scriptsize(55\%)},
                 {10--15\\\scriptsize(23\%)}, {$\geq$15\\\scriptsize(3\%)}, {All\\\scriptsize(100\%)}},
    xticklabel style={align=center, font=\footnotesize},
    enlarge x limits=0.15,
    axis lines*=left,
    nodes near coords={\pgfmathprintnumber{\pgfplotspointmeta}\%},
    nodes near coords style={font=\footnotesize},
    legend style={at={(0.5,1.03)}, anchor=south, legend columns=2, font=\scriptsize, draw=none},
]
\addplot[fill=gray!55, draw=black] coordinates {(1,20) (2,45) (3,64) (4,87) (5,47)};
\addlegendentry{drift within $\pm45^\circ$ of downwind}
\addlegendimage{line legend, dashed, black}
\addlegendentry{no wind effect (25\%)}
\draw[dashed] (rel axis cs:0,0.25) -- (rel axis cs:1,0.25);
\end{axis}
\end{tikzpicture}
\Description{Bar chart of the share of time the USV drifts within 45 degrees of downwind:
20\% at 1--5 kn, 45\% at 5--10 kn, 64\% at 10--15 kn, 87\% above 15 kn and 47\% overall,
against 25\% expected if wind had no effect on direction.}
\caption{Time the USV drifts within $\pm45^\circ$ of downwind, by wind speed
(share of minutes in brackets). Dashed: no wind effect.}
\label{fig:withwind}
\end{figure}

% NOTE: one figure: mean cross-validation error against the input window and against the learning rate,
% one line for each model. Say which setting is best, how flat the curve is, and why.
% [Claude] NOTE for Rowan: fig:hyper no longer exists, so this \ref prints "Figure ??" until the figure is added
% or the sentence changes. No sweep has been run yet (see the Hyperparameters note).
\textbf{RQ3: how sensitive are the models?} Figure~\ref{fig:hyper} shows the mean cross-validation error for each setting. 

{\color{red}
------------------}

{\color{red}
------------------}

{\color{red}
------------------}


\noindent\textbf{RQ4:} how does this differ from existing methods?
% This could be dropped because we dont have an experiemnt or proof or data

{\color{red}
------------------}

{\color{red}
------------------}

{\color{red}
------------------}


% NOTE: 20% of the mark. Give the one command that reproduces every table and figure, the fixed seeds,
% the run time, and where the results are written.
\textbf{Reproducibility.} The command \texttt{bash quickstart.sh} installs the pinned dependencies and reruns every step, from the raw telemetry in \texttt{data/} to the scores in \texttt{evaluation/results/} and the trained models in \texttt{output/}. Seeds are fixed, and the full run takes about 5~minutes on a laptop CPU.

% NOTE: if there is room, close with one sentence: what a pilot can rely on, and what comes next.

\bibliographystyle{ACM-Reference-Format}
\bibliography{bibliography.bib}

\end{document}
